\section{Introduction}
One of the technical barriers in deep learning theory is that the relationship between parameters and functions is a black box. For this reason, the majority of authors build their theories on extremely simplified mathematical models. Such theories can explain the complex phenomena in deep learning only at a highly abstract level.

The proof of a universality theorem contains hints for understanding the internal data processing mechanisms inside neural networks. For example, the first universality theorem\emph{s} for depth-2 neural networks were shown in 1989 with \emph{four} different proofs by \citet{Cybenko1989}, \citet{Hornik1989}, \citet{Funahashi1989}, and \citet{Carroll.Dickinson}. Among them, Cybenko and Hornik et al. presented existential proofs by using Hahn-Banach and Stone-Weierstrass respectively, meaning that it is not clear how to assign the parameters. On the other hand, Funahashi and Carroll-and-Dickinson presented constructive proofs by reducing networks to the Fourier transform and Radon transform respectively, meaning that it is clear how to assign the parameters. The latter constructive methods were refined as the so-called integral representation by \citet{Barron1993} and further culminated as the \emph{ridgelet transform}, the main objective of this study, discovered by \citet{Murata1996} and \citet{Candes.PhD}.

To show the universality in a constructive manner, we formulate the problem as a functional equation: 
Let $\lm[\gamma]$ denote a certain learning machine (such as a deep network) with parameter $\gamma$, and let $\calF$ denote a class of functions to be expressed by the learning machine. Given a function $f \in \calF$, find an unknown parameter $\gamma$ so that the machine $\lm[\gamma]$ represents function $f$, i.e.
\begin{align*}
	\lm[\gamma] = f,
\end{align*}
which we call a \emph{learning equation}. 
This equation is a stronger formulation of learning than an ordinary formulation such as minimizing empirical risk $\sum_{i=1}^n|\lm[\gamma](x_i) - f(x_i)|^2$ with respect to $\gamma$, as the latter is a weak form (or a variational form) of this equation. Therefore, characterizing the solution space of this equation leads to understanding the parameters obtained by risk minimization. Following previous studies \citep{Murata1996,Candes.PhD,Sonoda2021aistats,Sonoda2021ghost,Sonoda2022gconv,Sonoda2022symmetric}, we call a solution operator $\ridge$ satisfying $\lm[\ridge[f]] = f$ a \emph{ridgelet transform}.
Once such an $\ridge$ is found in a closed-form manner, we can present a constructive proof of \emph{universality} because the reconstruction formula $\lm[\ridge[f]]=f$ implies for any $f \in \calF$ there exists a machine that implements $f$. 

For depth-2 neural networks, the equation has been solved with several closed-form ridgelet transforms by using either Fourier expression method \citep{Sonoda2024fourier}, or group representation method \citep{sonoda2023joint}. For example, the closed-form ridgelet transforms have been obtained for depth-2 fully-connected networks \citep{Sonoda2021ghost}, depth-2 fully-connected networks on manifolds \citep{Sonoda2022symmetric}, depth-2 group convolution networks \citep{Sonoda2022gconv}, and depth-2 fully-connected networks on finite fields \citep{Yamasaki2023icml}.
Furthermore, \citet{Sonoda2021aistats} have revealed that the distribution of parameters inside depth-2 fully-connected networks obtained by empirical risk minimization asymptotically converges to the ridgelet transform. In other words, the ridgelet transform can also explain the solutions obtained by risk minimization.

On the other hand, for depth-$n$ neural networks, the equation is far from solved, and it is common to either consider infinitely-deep mathematical models such as Neural ODEs \citep{Sonoda2017otml,E2017,Li2018b,Haber2017,Chen2018a}, or handcraft networks that approximate another universal approximators such as piecewise polynomial functions and indicator functions.
For example, construction methods such as the Telgarsky sawtooth function (tent map, or the Yarotsky scheme) and bit extraction techniques
\citep{Cohen2016,Telgarsky2016,Yarotsky2017,Yarotsky2018,Yarotsky2020,Daubechies2022,Cohen2022,Siegel2023,Petrova2023,Grohs2023}
have been developed (not only to investigate the expressivity but also) to demonstrate the depth separation, super-convergence, and minmax optimality of deep ReLU networks. Various feature maps have also been handcrafted in the contexts of
geometric deep learning \citep{Bronstein2021gdl} and deep narrow networks \citep{Lu2017expressive,Hanin2017,Lin2018,Kidger2020,park2021minimum,Li2023minimum,cai2023achieve,kim2024minimum}. However, for the purpose of understanding the parameters obtained by risk minimization, these results are less satisfactory because there is no guarantee that these handcrafted solutions are obtained by risk minimization in a manner presented by \citet{Sonoda2021aistats}.


In order to investigate the relation between parameters and functions, we need to write down a general solution (i.e., the ridgelet transform) rather than handcrafting a particular solution. However, conventional ridgelet transforms have been limited to \emph{depth-2} networks. In other words, existing methods cannot construct solutions for networks that repeatedly compose nonlinear activation functions more than twice---such as $\sigma(A_2 \sigma(A_1 \xx - \bb_1) - \bb_2)$. In this study, inspired by the group-theoretic approach of \citet{sonoda2023joint}, we derive the ridgelet transform for \emph{depth-$n$} learning machines.


The contributions of this study are summarized as follows.
\begin{itemize}
    \item We derive the ridgelet transform (solution operator for learning equation) for a general class of learning machines called the \emph{joint-group-equivariant machine} (\cref{thm:main}), 
    which shows the universal approximation theorem for a wide range of learning machines in a constructive and unified manner.
    \item As applications, we show the universal approximation properties of \emph{four} examples: depth-$n$ joint-equivariant machine (\cref{sec:deep-ridgelet}), depth-$n$ fully-connected network (in \cref{sec:ex-depth-n}), depth-$n$ group-convolutional network (in \cref{sec:gconv}), and a new depth-$2$ network with quadratic forms whose universality has not been known (in \cref{sec:quad}).
\end{itemize}
    Until this study, the universality of deep networks has been shown in a different manner from the universality of shallow networks, but our results discuss them on common ground. Now we can understand the approximation schemes of various learning machines in a unified manner.




